\nonstopmode 
%\documentclass[11pt,a4paper,twoside]{tesis}
% SI NO PENSAS IMPRIMIRLO EN FORMATO LIBRO PODES USAR
\documentclass[11pt,a4paper]{tesis}
\usepackage{titlesec}
\titleformat{\chapter}[hang]{\normalfont\huge\bfseries}{\thechapter}{1em}{}
\titlespacing*{\chapter}{0pt}{0pt}{20pt}

\makeatletter
\renewcommand\chapter{\if@openright\else\fi
  \thispagestyle{plain}%
  \global\@topnum\z@
  \@afterindentfalse
  \secdef\@chapter\@schapter}
\makeatother

\usepackage{graphicx}
\usepackage[utf8]{inputenc}
\usepackage[spanish]{babel}
\usepackage[left=3cm,right=3cm,bottom=3.5cm,top=3.5cm]{geometry}

\begin{document}

%%%% CARATULA

\def\autor{Lautaro Lasorsa}
\def\tituloTesis{Propuesta de mejora en la medición del bienestar económico}
\def\runtitulo{Propuesta de mejora en la medición del bienestar económico}
\def\director{Rodrigo Castro}
\def\codirector{Walter Sosa Escudero}
\def\lugar{Buenos Aires, Argentina, 2024}
\newcommand{\HRule}{\rule{\linewidth}{0.2mm}}
%
\thispagestyle{empty}

\begin{center}\leavevmode

\vspace{-2cm}

\begin{tabular}{l}
\includegraphics[width=2.6cm]{logofcen.pdf}
\end{tabular}


{\large \sc Universidad de Buenos Aires

Facultad de Ciencias Exactas y Naturales}

\vspace{7.0cm}

%\vspace{3.0cm}
%{
%\Large \color{red}
%\begin{tabular}{|p{2cm}cp{2cm}|}
%\hline
%& Pre-Final Version: \today &\\
%\hline
%\end{tabular}
%}
%\vspace{2.5cm}

\begin{huge}
\textbf{\tituloTesis}
\end{huge}

\vspace{2cm}

{\large Tesis de Licenciatura en Ciencias de Datos}

\vspace{2cm}

{\Large \autor}

\end{center}

\vfill

{\large

{Director: \director}

\vspace{.2cm}

{Codirector: \codirector}

\vspace{.2cm}

\lugar
}

\newpage\thispagestyle{empty}


%%%% ABSTRACTS, AGRADECIMIENTOS Y DEDICATORIA
% \frontmatter

% \tableofcontents

% \mainmatter
% \pagestyle{headings}

%%%% ACA VA EL CONTENIDO DE LA TESIS

\chapter{Medición del Bienestar Ecónomico}

En este cápitulo se presenta la problematica a tratar en la tésis, se motiva su estudio y se comparán las formas actuales de medirlo con las alternativas propuestas.

\section{Definición y motivación}

En esta sección definimos qué es el bienestar económico y otros conceptos relacionados y por qué es importante medirlo bien.

\section{Medición actual}

En esta sección comentamos la forma actualmente utilizada, entre otros por quienes elaboran IDH, para medir el bienestar económico promedio de una población, analizando sus ventajas y desventajas.

\section{Propuesta de mejora: Medición ideal}

En esta sección comentaremos cuál es la que proponemos como la forma ideal de medir el bienestar económico, la justificación de esta propuesta y sus límitaciones, sobre todo a nivel logistico.

\section{Propuesta de mejora: Mediciones aproximadas}

En esta sección introducimos el concepto de granularidad de una medición, y estudiamos soluciones de compromiso para crear una medición que se aproxima más a la que proponemos como medición ideal pero siendo logisticamente factibles.
Serán las mediciones que utilizaremos con los datos empiricos.

\chapter{Datos sinteticos}

En este cápitulo se utilizán datos simulados para explorar caracteristicas de los índicadores (existentes y propuestos) en sí mismos.

\section{Generación}
En esta sección se comentará qué datos sinteticos se generaron, que asunciones se hicieron sobre la distribución de los mismos, y un breve comentario sober las tecnologías utilizadas.

\section{Datos $LN(0,1)$}

En esta sección se analizán los datos generados utilizando que los ingresos de una población siguen una distribución $LogNormal(0,1)$.

\section{Datos $LN(0,\sigma^2)$}

En esta sección se analizán los datos generados utilizando una distribución $LogNormal(0,\sigma^2)$ para distintos valores de $\sigma^2$. Trabajar con distintos valores de $\sigma^2$ permite una mayor variedad y riqueza de análisis.

\chapter{Datos reales: Distribución de los ingresos}

En esta sección se utilizán datos reales de distribución del ingreso de distintos países y años para estudiar su distribución y compararlos con los datos sinteticos antes generados.

\section{Datos disponibles}

Se presenta la base de datos disponible. Se comentan los distintos alcances (nacional, urbano y rural) y tipos (ingresos y consumo) de las encuestas disponibles.

\section{Distribución empirica}

Se estudian las distribuciones empiricas disponibles y se las compara con los datos sinteticos generados en el cápitulo anterior. Como el año es una variable importante para utilizar como variable de control, se estudiará la distribución condicional de los datos al año.

\chapter{Datos reales: Mejora de la capacidad predictiva}

En este cápitulo exploramos si las modificaciones propuestas en la medición del bienestar económico mejoran la capacidad de predicir otros indicadores socio económicos.

\section{Esperanza de Vida al Nacer (EVN)}

Se estudia pormenorizadamente el caso de la EVN, discutiendo distintas formas de evaluar la capacidad predictiva de las distintas mediciones. Se incluye el uso de bootstrap no parametrico para estimar la distribución de la correlación condicional (al año) entre las medidas con distinta granularidad y la EVN, siendo esta la metrica más madura y que se generalizará a otras variables socio economicas.

\section{Otros indicadores}

Utitizando la metrica comparativa finalmente desarrollada para el caso de la EVN, se aplica la misma metrica a otros indicadores socio economicos (ej: Alfabetización a los 15 años) para ver cómo se comportan es estos casos

\chapter{Conclusiones}

Titulo autoexplicativo

\chapter{Posibles trabajos futuros}

Titulo autoexplicativo

\chapter{Citas bibliograficas}

Titulo autoexplicativo



\end{document}
